\documentclass[11pt]{article}
\usepackage[margin=1in]{geometry}
\usepackage{amsmath,amssymb,bm}
\title{Magnetic field of the two-dimensional loop}
\author{}
\date{}
\begin{document}
\maketitle

\section{Problem Setup}
The conduction calculation gives an in-plane current
$\bm J=(J_x,J_y,0)$. The model is unchanged along the $z$
direction and assumes uniform permeability $\mu$, with no
applied magnetic field. Under these assumptions, the magnetic
field has only a $z$ component:
$\bm B=(0,0,B_z(x,y))$.

Ampere's law, $\nabla\times\bm B=\mu\bm J$, then becomes
\begin{equation}
 (\partial_yB_z,-\partial_xB_z)=\mu(J_x,J_y),
 \qquad
 \nabla B_z=\mu(-J_y,J_x).
 \label{eq:ampere}
\end{equation}
There is no current in air, so $\nabla B_z=0$ there: $B_z$
is constant in each connected air region. The outside constant
is set to zero. The cavity constant, called $C_B$, must be
determined from the current. With no idealized sheet current
at the conductor surfaces, $B_z$ is continuous. Thus it is
zero on the outer conductor surface and $C_B$ on the
cavity-facing surface.

\section{Fit the $B_z$ field}
Equation~\eqref{eq:ampere} suggests integrating
$\mu(-J_y,J_x)$ along a path to obtain $B_z$. However,
due to numeric errors of the finite-element method,
different paths could give different approximate values.
This motivates an optimization approach below.

Let $\bm g_{*,h}=\mu(-J_{y,h},J_{x,h})$ be the gradient
requested by Ampere's law in conductor triangle $h$.
Represent $B_z$ by a continuous piecewise-linear function
$\widetilde B_z = \sum_{i}b_iN_i$,
where the summation is over all conductor interior vertices
and conductor-cavity vertices. We further constraint all
conductor-cavity vertices share the same $b_i=C_B$.
With this choice of vertices, the boundary conditions are
automatically satisfied, that is,
\begin{equation}
 \widetilde B_z|_{\Gamma_{\rm outer}}=0,\quad
 \widetilde B_z|_{\Gamma_{\rm cavity}}=C_B.
\end{equation}

Our goal is to minimize the area-weighted
difference:
\begin{equation}
 \min_{\widetilde B_z}
 \frac12\sum_{h\in\mathcal T_c}A_h
 \left|\sum_{i=0}^2 b_{h,i}\nabla N_{h,i}-\bm g_{*,h}\right|^2,
 \label{eq:fit}
\end{equation}
Here $\mathcal T_c$ is the set of conductor triangles, and $b_{h,i}$
is the value of $\widetilde B_z$ at local vertex $i$ of triangle $h$.
It is understood that $b_{h,i}=0$ if the vertex falls onto conductor-outer
boundary and $b_{h,i}=C_B$ if the vertex falls onto conductor-cavity boundary.
Taking derivative to Equation \ref{eq:fit} gives a linear system that
we can solve to get $\widetilde B_z$.

\section{Field used by the app}
The app stores the solved $B_z$ at conductor vertices, fills
cavity vertices with $C_B$, and fills outside vertices with
zero. The magnetic view interpolates those nodal values
within each triangle. The energy-flow calculation uses the
triangle average
\begin{equation}
 \overline B_{z,h}=\frac{b_{h,0}+b_{h,1}+b_{h,2}}{3}.
\end{equation}
This is a strictly two-dimensional magnetic model; it is not
the magnetic field of a finite three-dimensional wire.

\paragraph{Code:} \texttt{calculations/magnetics.py}.
\end{document}
